\chapter{迁移学习：复用表示，再验证目标任务}
\label{ch:transfer}

迁移学习（Transfer Learning）把已有任务或数据中学到的信息用于目标任务。
常见做法是使用预训练特征提取器，接上适合新标签的任务头；
不要求源与目标的类别名称相同，但输入与表示需要适配。

\section{冻结与微调分别更新什么？}
\label{sec:transfer-freeze}
\begin{figure}[H]
\centering\begin{tikzpicture}[node distance=12mm]
\node[box,text width=25mm] (x) {目标任务输入};
\node[box,right=of x,text width=40mm] (base) {预训练特征提取器\\输出$d$维表示};
\node[box,right=of base,text width=32mm] (head) {新任务头\\输出$K$类logits};
\draw[flow] (x)--(base);\draw[flow] (base)--(head);
\node[lab,below=10mm of base,text width=68mm,align=center] {冻结：不更新底座参数；微调：更新选定参数。\\两种情况都要规定训练/评估模式与预处理。};
\end{tikzpicture}

\caption{目标任务有$K$个互斥类别时，线性任务头将$d$维表示映成$K$个logits。}
\end{figure}
\readfigure{若$d=16,K=3$，带偏置任务头有$16\times3+3=51$个参数。
冻结底座时只训练这些新参数；部分或全部微调时，还更新被选中的底座参数。
冻结减少可训练参数，但不消除底座前向计算。}

\begin{center}
\begin{tabularx}{\textwidth}{P{31mm}Y}
\toprule
策略 & 需要比较的因素\\
\midrule
冻结表示，仅训练头 & 可作较省训练计算的起点，表示与目标是否匹配仍需验证。\\
解冻部分层 & 控制适应程度、学习率与训练模式；没有通用解冻层数。\\
全部微调 & 更大适应空间，也需要数据、计算与过拟合检查。\\
\bottomrule
\end{tabularx}
\end{center}

冻结参数与评估模式是两回事：BatchNorm运行统计是否更新、Dropout是否随机，
需要另外设定；不能只设置“不求参数梯度”就默认底座输出固定。
预处理、词表/分词和输入尺度也需与预训练约定衔接。

\begin{pitfall}[预训练不保证每个目标任务都有收益]
源与目标差异、标签质量和数据量会影响结果，可能出现负迁移。
应与简单基线比较，按目标对象或时间划分验证，仍要隔离测试集。
还需核对预训练数据是否与评估样本重叠，避免把重复内容带来的好分数当作泛化。
\end{pitfall}

\noindent\textbf{参考：}《动手学深度学习》微调章节；PyTorch Transfer Learning Tutorial。

